\documentclass[12pt,a4paper]{article}
\usepackage{iftex}
\ifPDFTeX
    \usepackage[utf8]{inputenc}
    \usepackage[T1]{fontenc}
    \usepackage{mathpazo}
\else
    \usepackage{fontspec}
    \setmainfont{Times New Roman}
\fi
\usepackage[brazilian]{babel}
\usepackage{amsmath, amssymb, amsthm}
\usepackage[linesnumbered,ruled,vlined,portuguese]{algorithm2e}
\usepackage{listings}
\usepackage{booktabs}
\usepackage{geometry}
\usepackage{float}
\usepackage{xcolor}
\usepackage{setspace}
\usepackage[colorlinks=true, linkcolor=blue, urlcolor=blue, citecolor=blue]{hyperref}

\geometry{margin=2.5cm}
\onehalfspacing

\definecolor{codebg}{rgb}{0.96,0.96,0.96}
\definecolor{codekey}{rgb}{0.0,0.4,0.8}
\definecolor{codestr}{rgb}{0.6,0.1,0.1}
\definecolor{codecom}{rgb}{0.4,0.6,0.4}

\lstset{
    backgroundcolor=\color{codebg},
    commentstyle=\itshape\color{codecom},
    keywordstyle=\bfseries\color{codekey},
    stringstyle=\color{codestr},
    basicstyle=\ttfamily\footnotesize,
    breaklines=true,
    numbers=left,
    numberstyle=\tiny\color{gray},
    frame=single,
    rulecolor=\color{lightgray},
    captionpos=b,
    tabsize=4
}

\title{
    \vspace{-1cm}
    \begin{center}
        \Large \textbf{Universidade Federal do Ceará -- UFC} \\
        \large Departamento de Computação \\
        \vspace{1cm}
        \Huge \textbf{Métodos Numéricos II} \\
        \vspace{0.5cm}
        \Large \textit{Relatório da Unidade 4: Problemas de Valor Inicial (PVI)}
    \end{center}
}
\author{
    Nicolas Wallace Amorim Perote -- Matrícula: 566512 \\
    Samyra Vitória Lima de Almeida -- Matrícula: 521240
}
\date{\today}

\begin{document}

\maketitle
\tableofcontents
\newpage

\section{Introdução}
O estudo numérico de Problemas de Valor Inicial (PVI) é imprescindível para prever a evolução determinística de sistemas regidos por Equações Diferenciais Ordinárias (EDOs). Formalmente, buscamos projetar a curva $y(t)$ que obedece a um gradiente $y'(t)=f(t,y)$, originada num estado absoluto conhecido $y(t_0)=y_0$.

Nesta \textbf{Unidade 4}, desenvolvemos de forma unificada uma biblioteca de propagações temporais (A classe \texttt{IVPSolvers}), integrando métodos explícitos (Euler Clássico e Runge-Kutta 3/4) e métodos implícitos com reciclo histórico (Euler Implícito e Preditor-Corretor). A arquitetura vetorizada em \textit{NumPy} garante escalabilidade irrestrita para resolver desde problemas escalares simples até complexos arranjos multifásicos acoplados (Sistemas EDO).

\section{Métodos de Euler}
\subsection{Abordagem Explícita}
O Método de Euler baseia-se num truncamento raso da Série de Taylor, absorvendo tão-somente a derivada de primeira ordem para guiar a projeção linear iterativa.
\begin{equation}
y_{k+1}=y_k+\Delta t\,f(t_k,y_k)
\end{equation}
Seu baixo custo computacional (uma única avaliação do gradiente $f$ por passo) o torna atrativo para prototipagem rápida; contudo, a dependência geométrica estrita da tangente imediata propaga erro global de ordem $\mathcal{O}(\Delta t)$.

\subsection{Abordagem Implícita}
Para amortecer oscilações espúrias em sistemas altamente rígidos (\textit{stiff}), o Euler Implícito afere o declive no \textit{futuro} instante de avaliação.
\begin{equation}
y_{k+1}=y_k+\Delta t\,f(t_{k+1},y_{k+1})
\end{equation}
Isso levanta a necessidade de solucionar a incógnita isoladamente em ambos os flancos. Para EDOs não-lineares, arquitetou-se na nossa biblioteca o solucionador baseado no Teorema de Ponto Fixo: usando o limite explícito como palpite de inicialização $y_{k+1}^{(0)} = y_k$, efetuamos o arrasto:
\begin{equation}
y_{k+1}^{(i+1)} = y_k + \Delta t\, f(t_{k+1}, y_{k+1}^{(i)})
\end{equation}
até que a variação resulte diminuta (abaixo da tolerância \texttt{tol=1e-8}).

\section{Família Runge-Kutta}
Visando expandir o erro global mitigado e a precisão da aproximação em detrimento da escassa linearidade, as abordagens de Carl Runge e Martin Kutta ponderam múltiplas derivadas em subespaços do passo original.

\subsection{Runge-Kutta de 3ª Ordem (RK3)}
O RK3 colapsa a incerteza utilizando uma amostragem tripartite. Em nossa execução vetorizada (com matrizes \texttt{NumPy}), os ganhos avaliados $k_1, k_2, k_3$ modelam uma parábola local cujo peso atualiza o sistema:
\begin{equation}
y_{k+1}=y_k+\frac{\Delta t}{6}(k_1+4k_2+k_3)
\end{equation}

\subsection{Runge-Kutta de 4ª Ordem (RK4)}
O ápice da acurácia passo-a-passo no nosso escopo é detido pelo vetorizado RK4. Ele interroga o sistema em 4 instantes (início, duas avaliações centrais e término) com pesos alinhados aos preceitos da integração analítica de Simpson. 
\begin{equation}
y_{k+1}=y_k+\frac{\Delta t}{6}(k_1+2k_2+2k_3+k_4)
\end{equation}
Empregamos este algoritmo como lastro para incrustar as fases preparatórias nos métodos implícitos subsequentes, provendo margens de erro formidáveis $\mathcal{O}(\Delta t^4)$.

\section{Algoritmo Preditor-Corretor (PC4)}
Os métodos de Adams-Bashforth e Moulton pertencem à classe Multistep, baseados na integração formal de polinômios interpoladores retroativos. O método \textbf{PC4} tira benefício de computações do passado para projetar as retas.

\begin{enumerate}
    \item \textbf{(Preditor)} Com Adams-Bashforth de quatro passos, prevemos $\hat{y}_{k+1}$ avaliando explicitamente as 4 etapas prévias com coeficientes estritos $[55, -59, 37, -9]$.
    \item \textbf{(Corretor)} Empregando Adams-Moulton (implícito), iteramos o cômputo da aproximação $\hat{y}_{k+1}$ com um refinamento iterativo dependente do próprio gradiente futuro e de 3 etapas prévias $[9, 19, -5, 1]$.
\end{enumerate}

Para impulsionar a engrenagem (Arrancada), implementamos a resolução dos primeiros 4 estados autônomos exclusivamente via \textbf{RK4}.

\section{Experimentações e Desempenho}

\subsection{PVI-1: Equação Diferencial Escalar}
Avaliamos a equação $y' = \frac{2}{3}y$ partindo de $y(0) = 2$.
Com passo temporal agressivo $\Delta t = 0.1$, a solução teórica requer convergence a $414.254498$ no tempo 8.

\begin{itemize}
    \item \textbf{Euler Explícito}: Convergiu em $\approx 349.40$ (Erro catastrófico de 64 pontos).
    \item \textbf{Euler Implícito}: Fechou em $\approx 498.98$ (Desvio expressivo).
    \item \textbf{RK4}: Cravou impressionantes $\approx 414.254154$ (Erro marginal de $3.44 \times 10^{-4}$).
    \item \textbf{Pred.-Corretor (PC4)}: Aderiu fortemente em $\approx 414.255537$ (Erro de $1.04 \times 10^{-3}$).
\end{itemize}
Com $\Delta t = 0.01$, RK4 e PC4 atingiram margens de erro residuais de notáveis $10^{-7}$.

\subsection{PVI-2: Dinâmica e Balística Acoplada}
O método é provado na prática mapeando um vetor acoplado $\vec{S} = [v, y]$, modelando um projétil restrito por resistência atmosférica sob constante elástica aerodinâmica, descrito pelo subsistema EDO acoplado de arrasto $\mathcal{F}(\vec{S})$.

Em $\Delta t = 0.001$:
\begin{itemize}
    \item \textbf{RK3 e PC4} concordaram integralmente nas observações: alcançaram a altura de crista $y_{max} \approx 201.200242$ m.
    \item Na intersecção com o chão, ambos reportaram tempo de colisão iterado exato $t_{imp} \approx 7.789758$ s.
    \item A velocidade terminal se enrijeceu em $v_{imp} \approx -47.897575$ m/s.
\end{itemize}

\section{Conclusão}
O arcabouço computacional da Unidade 4 sintetiza a transição da álgebra simples para predição estocástica robusta. O método Euler pavimenta o fluxo lógico mas desidrata-se perante severas curvas transientes. O Runge-Kutta de quarta-ordem consagra o equilíbrio imutável entre complexidade per capita de processamento e extrema adesão de resposta, enquanto arranjos como Preditor-Corretor viabilizam o tráfego ótimo com enorme economia de ciclos do gradiente. A modelagem orientada a matrizes proporcionou uma biblioteca irretocável adaptada a sistemas multifísicos irrestritos.

\end{document}
